\documentclass[12pt]{article}
\usepackage[utf8]{inputenc}
\usepackage[T1]{fontenc}
\usepackage{hyperref}
\title{Wayu Language Resources for K-12 Curriculum Development}
\author{Compiled by Claude Sonnet 5 (claude-sonnet-5) for Andr\'as Kornai}
\date{}
\begin{document}
\maketitle

\section{Introduction}
Wayu (also spelled Vayu or Hayu; ISO 639-3 vay) is a Sino-Tibetan language of Nepal, written in Devanagari. Sources found in search give about 1,500 speakers in the 2011 census, well below the figure of about 10,000 in the survey brief; the discrepancy is not resolved here. One source reports that since 1973 most children in Murajor no longer speak it, while it is still being learnt in Manedihi. Do not confuse it with Raute, a different language of Nepal. Search found no school materials, and little else, in Wayu.

\section{Primary Educational Resources}
No primer, textbook or school programme for Wayu was found in search results.

\section{Linguistic and Reference Materials}
\begin{itemize}
\item Wikipedia, ``Hayu'': general overview, which cites the descriptive work of Boyd Michailovsky (including the 1988 book \emph{La langue hayu} and the Hayu chapter in \emph{The Sino-Tibetan Languages}, Routledge 2003). \url{https://en.wikipedia.org/wiki/Hayu}
\item Wikipedia, ``Vayu language'': overview of the language. \url{https://en.wikipedia.org/wiki/Vayu_language}
\item Max Planck EVA, numeral system note for Wayu. \url{https://lingweb.eva.mpg.de/channumerals/Wayu.htm}
\end{itemize}
The grammar and dictionary works cited in these pages were not themselves found online in search, so no direct links are given.

\section{Religious and Community Resources}
No scripture translation or community organisation specific to Wayu was found in search results.

\section{Digital and Audio Resources}
\begin{itemize}
\item Pangloss Collection (LACITO-CNRS): digitised Michailovsky field recording, ``Hayu 1972 Tape 2 Side A (elicitation)''. \url{https://cocoon.huma-num.fr/exist/crdo/meta/cocoon-7833b117-00f3-3885-a7aa-23c32f6d250e}
\item Endangered Languages Project entry for Hayu, including samples. \url{https://www.endangeredlanguages.com/lang/3399}
\end{itemize}

\section{Curriculum Development Context}
The language is endangered and is not known to be taught in schools. Documentation is mostly linguistic (Michailovsky) and archival. A curriculum would depend on new community-led materials, built with speakers in Murajor and Manedihi and starting from the existing archival recordings and published descriptions.

\end{document}